\section{Proposta de solução}

\subsection{Arquitetura do chatbot legislativo}

A solução proposta materializa-se em um chatbot legislativo que utiliza a abordagem RAG para responder a consultas sobre discursos da 56\textordfeminine{} Legislatura do Senado Federal. A arquitetura compreende três serviços principais: uma instância do Open WebUI, utilizada como interface de interação e orquestração do fluxo de RAG; scripts Python empregados na extração e preparação de conteúdo documental; o ChromaDB, utilizado como banco vetorial acessado por HTTP; o Ollama para integração com grandes modelos de linguagem; e a API da OpenAI tanto para geração dos embeddings quanto para geração de respostas.

O ambiente foi definido para rodar em ambiente local usando Docker e preparado para operar tanto com embeddings locais, via Ollama, quanto com embeddings providos por serviços externos, como a API da OpenAI. Na configuração experimental validada, adotou-se esta segunda opção. Essa flexibilidade amplia a adaptabilidade da solução a diferentes restrições institucionais de infraestrutura, custo ou política tecnológica.

\subsection{Pipeline desenvolvido}

O pipeline do projeto inicia-se com a aquisição dos dados, baseada no dataset \nolinkurl{fabriciosantana/discursos-senado-legislatura-56}, que reúne discursos da 56\textordfeminine{} Legislatura do Senado Federal. Salienta-se que esse dataset foi construído em trabalho anterior a partir do portal de dados abertos do Senado. No contexto deste projeto, na etapa de pré-processamento e seleção textual, priorizou-se o campo ``TextoDiscursoIntegral'', recorrendo-se aos campos ``Resumo'' e ``Indexacao'' quando necessário, de modo a privilegiar o conteúdo mais completo e informativo na construção da base.

Em seguida, realizou-se o chunking, com geração de segmentos que preservam metadados relevantes por chunk. Esses metadados incluem data, autor, partido, unidade da federação, tipo de uso da palavra, resumo, indexação e URL do texto integral. Tal modelagem favorece não apenas a recuperação semântica, mas também a auditabilidade das respostas, uma vez que permite reconstituir a origem documental dos trechos utilizados.

A partir desses segmentos, foram gerados arquivos de ingestão para o Open WebUI em lotes estruturados em formato Markdown, além de um arquivo \texttt{jsonl} de referência. A importação dos lotes foi automatizada por script responsável por realizar o upload, monitorar o processamento e garantir a inclusão dos conteúdos na \textit{knowledge base} do Open WebUI. Concluída a ingestão, procedeu-se à indexação vetorial, com geração e persistência dos embeddings no ChromaDB.

Na etapa de consulta, o fluxo RAG integra recuperação vetorial, montagem de contexto, aplicação de prompt e geração da resposta pelo modelo. Por fim, a avaliação inicial foi realizada por meio do script \texttt{run\_rag\_eval.py}, que consulta a \textit{knowledge base} via API e produz saídas nos formatos \texttt{.jsonl}, \texttt{.md} e \texttt{.csv}, viabilizando análise estruturada, revisão qualitativa e pontuação manual para comparação entre execuções.

\subsection{Principais contribuições}

O projeto apresenta contribuições em dois níveis complementares: técnico e institucional. No plano técnico, constrói-se um ambiente reprodutível de chatbot legislativo com RAG adaptado a um acervo real do Senado Federal; demonstra-se a importância de uma modelagem de chunks enriquecida por metadados; automatizam-se etapas de importação, ingestão e avaliação; e evidencia-se empiricamente a sensibilidade do desempenho às escolhas de segmentação e conectividade com banco vetorial remoto.

No plano institucional, a solução demonstra a viabilidade de aplicar técnicas contemporâneas de recuperação e geração ao contexto do Parlamento brasileiro, aproximando recomendações gerais da literatura de RAG em produção da realidade de um acervo legislativo público. Desse modo, contribui-se para a discussão sobre governança, auditabilidade e uso responsável de inteligência artificial em contextos parlamentares.

\subsection{Riscos e limitações}

Os principais riscos e limitações identificados relacionam-se à sensibilidade do desempenho à estratégia de chunking, à possibilidade de combinação excessiva de fontes em perguntas muito abertas, à dependência da qualidade dos documentos e metadados de origem e ao risco de respostas excessivamente amplas ou inferenciais quando o prompt não é suficientemente restritivo. Em caso de limitação de acesso a recursos computacionais, há, ainda, dependência operacional de serviços externos, tanto para geração de embeddings quanto para parte da geração de respostas.

Esses riscos devem ser compreendidos em perspectiva ampliada. Como observam Deene (2025) e Said (2025), sistemas RAG podem melhorar controle e auditabilidade, mas não eliminam automaticamente problemas de qualidade documental, governança de dados, conformidade, monitoramento ou avaliação. Em outras palavras, o RAG reduz determinadas fragilidades do uso puro de LLMs, mas não substitui a necessidade de processos institucionais de revisão e controle.
